\section{A Fermat-only diagonal refinement}
\label{app:fermat-refinement}

This appendix records a stronger statement that uses more symmetry on the
Fermat rail than is common to the cubic--quadric pair.  It is not used in
the proof of Theorem~\ref{thm:denominator} or
Corollary~\ref{cor:no-central-sector}.

For the Fermat cubic in six variables, the projective diagonal group is
\[
D=(\mu_3^6)/\mu_3.
\]
The primitive character labels have
\[
(a_0,\ldots,a_5)\in\{1,2\}^6,\qquad
\sum_i a_i\equiv0\pmod3.
\]
Br\"unjes identifies the primitive character labels
\cite[Lemma/Definition 4.3]{Brunjes2003}.  The corresponding primitive
pieces are one-dimensional by
\cite[Theorem 4.6]{Brunjes2003}.  In this six-variable cubic, the
admissible tuples have respectively zero, three, or six entries equal to
\(2\), accounting for
\[
1+\binom63+1=22
\]
primitive lines, in agreement with the squarefree degrees \(0,3,6\) used
in Section~\ref{sec:n3-character}.

\begin{proposition}[Diagonal-sector obstruction]
\label{prop:fermat-diagonal}
After extending coefficients to a splitting field for \(D\), no nonzero
central sum of primitive Fermat diagonal-character sectors admits ordinary
multiplicity \(4/3\).
\end{proposition}

\begin{proof}
Every selected primitive character occurs with multiplicity one.  An actual
representation scaled by \(4/3\) would require its multiplicity to be
divisible by three, which fails on every selected line.  Packaging Galois
orbits over a smaller coefficient field does not help: after scalar
extension, each constituent again has multiplicity one.
\end{proof}

The classical interpretation of the associated Fermat eigenvalues by Jacobi
sums is described in \cite{Weil1952,Brunjes2003}, but no Jacobi-sum
analytic consequence is used here.  In particular, the proposition proves
neither automorphy nor a functional equation.  It is deliberately confined
to the Fermat diagonal group: the complete-intersection rail is stable only
under the common group \(G_3\), so the main two-rail obstruction must use
Theorem~\ref{thm:n3-character}.
